\subsection{An endpoint obstruction to global analytic decay}
\label{ssec:endpoint}

The logarithmic endpoint term imposes a restriction on possible coefficient
laws. It is stronger than the observation that finite-window fitted exponents
drift, and does not require an eigenfunction to have one sign.

\begin{proposition}[Smooth eigenfunctions must be flat at the boundary]
\label{prop:endpointflat}
Let $c>1$. If an eigenfunction $u$ of the exact full CCM operator has a
$C^\infty$ representative on $[0,L]$, $L=\log c$, then
$u^{(m)}(0)=u^{(m)}(L)=0$ for every $m\ge0$.
Consequently no nonzero eigenfunction is analytic on a neighborhood of
$[0,L]$, and no nonzero eigenfunction has periodic Fourier coefficients
bounded by $Ce^{-a|j|^q}$ for any $a>0$, $q\ge1$.
\end{proposition}
\begin{proof}
For the zero-extension realization $Q_a$ of
$m_a(t)=\tfrac12\log(a^2+t^2)$, the symmetric translation formula gives
\begin{align*}
 Q_a u(x)={}&(\log a)u(x)
 +\frac12\int_0^L e^{-a|x-y|}\frac{u(x)-u(y)}{|x-y|}\,dy\\
 &+\frac12u(x)\bigl(E_1(ax)+E_1(a(L-x))\bigr),\qquad 0<x<L.
\end{align*}
The two exponential integrals are exactly the contributions from outside
the interval. By the expansion of $E_1$ at zero
\cite[eq.~6.6.2]{DLMF}, $E_1(ax)+\log x$ is analytic near zero.
Split the integral at $y=x$; the divided difference extends smoothly
through the diagonal, since
\[
 \frac{u(x)-u(y)}{x-y}=\int_0^1u'(y+s(x-y))\,ds.
\]
It follows that $Q_au(x)=-\tfrac12u(x)\log x+G_a(x)$ with $G_a$
smooth up to the left endpoint whenever $u$ is smooth on $[0,L]$.

For $a=1/2$, the remaining archimedean kernel in
Proposition~\ref{prop:logdiag} is explicitly
\[
 -b(2|x-y|)e^{-|x-y|/2},\qquad
 b(s)=\frac1{1-e^{-s}}-\frac1s.
\]
This follows by Fourier inversion of the digamma integral used there.
The singularity of $b$ at zero is removable, with
$b(s)=\tfrac12+s/12+O(s^3)$. Splitting at $x=y$ again shows that this
remainder maps smooth functions to functions smooth at the endpoint.
The pole kernel is analytic. In a sufficiently small one-sided neighborhood
of zero, each of the finitely many prime translations is either zero or
an interior translate of $u$. A translation of length $L$ is zero on the
open interval. Thus the full eigenvalue equation is
\[
 0=(A_c-\varepsilon)u(x)=-\tfrac12u(x)\log x+G(x),
 \qquad G\in C^\infty([0,\eta])
\]
for some $\eta>0$. If $m$ is the first nonvanishing derivative of $u$ at
zero, the $m$th derivative of $u(x)\log x$ contains
$u^{(m)}(0)\log x$, which is unbounded. This contradicts the smoothness
of $G$. Hence $u$ is flat at zero, and reflection gives flatness at $L$.
An analytic function flat at an endpoint is identically zero on the
connected interval. A Fourier bound with $q=1$ gives a holomorphic strip
extension, while $q>1$ gives a holomorphic extension to every finite strip;
both are therefore excluded.
\end{proof}

This does not rule out $0<q<1$: smooth nonanalytic functions can be flat at
both endpoints. It also does not prove the nonzero logarithmic-boundary
amplitude needed for the conditional $1/[N(\log N)^2]$ rate above.

\begin{corollary}[No exact termination of the CCM Ritz sequence]
\label{cor:notermination}
For every finite $N$, the smallest exact Ritz value satisfies
$\varepsilon_N>\varepsilon_\infty$.
\end{corollary}
\begin{proof}
Equality in the variational principle would put a nonzero trigonometric
polynomial in the continuum ground eigenspace. Such a polynomial is
analytic, contradicting Proposition~\ref{prop:endpointflat}.
\end{proof}

The calibrated parameter $q_\gamma(c)=2/(2\gamma+\log\log c)$ satisfies
$q_\gamma(c)\ge1$ up to
$c_*=\exp(\exp(2-2\gamma))=10.27077047\ldots$, wherever its denominator
is positive. In particular $q_\gamma(7)=1.0988038\ldots$ and
$q_\gamma(9)=1.0300643\ldots$. Its proposed identification with a global
continuum coefficient exponent is therefore false at those cutoffs.
The finite-window root-error fit and its measured accuracy are separate:
a finite band can decay much faster than the eventual endpoint tail.
